% Build: pdflatex teaching-challenge-white-paper.tex  (run twice for references)
% Replace the author line before submitting.
\documentclass[10pt,letterpaper]{article}

\usepackage[left=0.85in,right=0.85in,top=0.7in,bottom=0.7in]{geometry}
\usepackage{booktabs}
\usepackage{tabularx}
\usepackage{enumitem}
\usepackage{titlesec}
\usepackage{multicol}
\usepackage{microtype}
\usepackage[hidelinks]{hyperref}

\titlespacing*{\section}{0pt}{1.1ex plus 0.3ex}{0.5ex}
\titlespacing*{\subsection}{0pt}{0.9ex plus 0.2ex}{0.4ex}
\titleformat*{\section}{\large\bfseries}
\titleformat*{\subsection}{\normalsize\bfseries}
\setlist[itemize]{topsep=0.3ex,itemsep=0.2ex,parsep=0pt,leftmargin=1.4em}
\setlength{\parskip}{0.4ex}
\renewcommand{\arraystretch}{1.05}
\setlength{\textfloatsep}{6pt}
\setlength{\intextsep}{6pt}
\setlength{\abovecaptionskip}{3pt}
\setlength{\belowcaptionskip}{2pt}
\newcolumntype{Y}{>{\raggedright\arraybackslash}X}

\title{\vspace{-3.5ex}\textbf{Industry experience infusion in graduate education:\\ why, how, when, and where}}
\author{\normalsize Group N: Member One, Member Two, Member Three\\[-0.4ex]
        \small George W. Woodruff School of Mechanical Engineering, Georgia Institute of Technology}
\date{\vspace{-1.5ex}\small\today}

\begin{document}
\maketitle
\vspace{-3ex}

\noindent\textit{Graduate engineering programs certify what students know and, for doctoral students,
that they can produce new knowledge. Neither credential says much about whether a graduate can take an
underspecified problem from a paying stakeholder and return a defensible answer on a deadline. We argue
the fix is not another semester-long course but a supply of short, mentor-reviewed industry projects
whose assessment produces reusable evidence, and we give an implementation strategy, an MS versus PhD
comparison, a Woodruff School adoption path, and a working case study.}

\section{Why: the gap a graduate curriculum leaves open}

Professional judgment is learned by doing bounded real work under someone who already has it, then
being told precisely where the work fell short. Graduate coursework rarely supplies that loop.
Problems arrive pre-scoped, the stakeholder is the instructor, and the feedback arrives once, as a
grade, after the work can no longer be revised. Experiential and situated accounts of learning predict
the outcome directly. A cycle that never closes through action and reflection, inside the community
that uses the skill, transfers poorly~\cite{kolb1984,lave1991,collins1989}.

The National Academies' review of graduate STEM education makes the same point, recommending that
programs take explicit responsibility for professional competencies instead of treating them as a
byproduct of research~\cite{nas2018}. Most engineering PhD graduates work outside the professoriate, so
preparation for an academic career is the minority case being optimized. For MS students the gap is
sharper still, since many enter a terminal degree specifically to change what they can be hired to do,
with three or four semesters to do it.

\section{How, part one: what is already being tried}

Industry-focused graduate learning is not a new idea, and the existing forms fail in instructive ways.
Table~\ref{tab:approaches} summarizes what is in use.

\begin{table}[ht]
\centering\scriptsize
\caption{Current approaches to industry-focused learning in graduate education.}
\label{tab:approaches}
\begin{tabularx}{\textwidth}{@{}>{\raggedright\arraybackslash}p{0.27\textwidth} Y Y@{}}
\toprule
\textbf{Approach} & \textbf{Pedagogical strength} & \textbf{Why it resists scale} \\
\midrule
Industry-sponsored capstone or practicum
 & Authentic task, real client, team coordination~\cite{crawley2014}
 & One sponsor per team per term; faculty effort scales with enrollment \\
Co-op and internship
 & Full immersion in a community of practice
 & Calendar-sized; competitive; learning rarely assessed \\
Industrial doctorate or sponsored fellowship (EngD, Mitacs, NSF INTERN)
 & Long horizon, genuine co-supervision
 & Negotiated one student at a time; IP terms take months \\
Professional MS practicum
 & Aligned with a terminal degree's purpose
 & Exists only where a program was built around it \\
Case-based or guest-lecture seminar
 & Cheap, scales to any enrollment
 & Students analyze work instead of doing it; no stakeholder \\
Consulting-style clinic
 & Repeat clients, cumulative reputation
 & Needs dedicated staff and a standing client pipeline \\
\bottomrule
\end{tabularx}
\end{table}

One pattern runs through all of it. Every approach with strong pedagogy is priced per student and
negotiated by a human, and every approach that scales cheaply has removed the stakeholder. The binding
constraint is not student interest or company goodwill. It is the coordination cost of matching,
supervising, and assessing one student against one real problem.

\section{How, part two: our proposal}

Unbundle ``industry experience'' from the semester. Instead of one large engagement per student per
term, build a supply of small projects, each with a published scope, a named industry mentor, a
scoring rubric tied to named skills, and an effort band between roughly 2 and 15 hours. Students
complete several over a degree. Assessment of each one produces a durable record of which skills were
demonstrated, on whose authority, against which criteria.

This is cognitive apprenticeship with the scaffolding made explicit and the assessment made
reusable~\cite{collins1989,wood1976}. The rubric is what makes it work at scale rather than the
software. Publishing criteria before the work starts is constructive alignment in its plainest
form~\cite{biggs1996}, it makes mentor feedback specific enough to act on, and it gives the program
something defensible to report~\cite{gulikers2004,jonsson2007}.

Five design commitments do most of the pedagogical work, and each one exists to block a specific
failure we expect:

\begin{itemize}
  \item \textbf{Evidence, not hours.} Credit attaches to a rubric score from a named assessor, never to
        logged time, which would reward slow work and unreliable self-reports.
  \item \textbf{One revision round is in scope by default.} Feedback a student cannot act on is a grade
        wearing a costume. Revision is where the learning happens~\cite{ericsson1993}.
  \item \textbf{Progression signals never gate access.} Points and levels may motivate, but only skill
        evidence and stated prerequisites decide what a student attempts next. Otherwise the credential
        becomes a measure of platform activity.
  \item \textbf{Two independent assessors for a durable claim.} One favorable review is provisional;
        confirmation needs a second project and a second assessor. This limits lenient mentors and
        lucky matches alike.
  \item \textbf{A named mentor and a named escalation owner on every engagement.} Projects stall for
        ordinary reasons, and deciding in advance who rescopes or reassigns keeps a stall from becoming
        the student's fault.
\end{itemize}

\subsection{Implementation strategy in a large program}

A large program cannot negotiate each placement, so five steps that currently consume faculty time have
to become routine operations. A coordinator recruits briefs and rejects any that lack a deliverable, a
rubric, and a named mentor with confirmed capacity. Students see why they are eligible and what
evidence a project would produce, so self-selection replaces assignment. Mentors score the published
rubric, with actionable feedback within five business days tracked as a program metric rather than a
hope. A standing recovery process extends, rescopes, reassigns, or closes an engagement without
penalizing the student. Recognition starts as a transcript notation or badge and converts to credit
only after a term of data exists, because asking a curriculum committee to approve credit for an
unproven mechanism is how good pilots die.

The honest cost is that coordinator. This model moves effort off faculty and onto operations, and a
program unwilling to fund roughly half a staff position should not start.

\section{When and where: MS versus PhD}

The mechanism is the same for both degrees; the reason for using it is not. Table~\ref{tab:msphd}
compares them.

\begin{table}[ht]
\centering\scriptsize
\caption{Utility of short mentor-reviewed industry projects by degree objective.}
\label{tab:msphd}
\begin{tabularx}{\textwidth}{@{}l Y Y@{}}
\toprule
 & \textbf{MS (coursework or terminal)} & \textbf{PhD} \\
\midrule
Degree objective
 & Advanced applied competence, then employment
 & Independent original research and a defended contribution \\
What the projects add
 & The core professional outcome the degree promises but rarely assesses
 & Breadth, stakeholder translation, and a credible non-academic option \\
Best timing
 & Term one onward; deliberately before the job search
 & After the qualifying exam and proposal, when the research plan can absorb interruption \\
Fit with existing structure
 & High; fills unstructured elective and professional-development space
 & Low; competes directly with advisor-funded research time \\
Effort band that works
 & Several small projects, stacking evidence across skills
 & Fewer, larger, deliberately adjacent to the dissertation topic \\
Main risk
 & Becomes r\'{e}sum\'{e} decoration if rubrics are weak
 & Advisor reads it as lost research months \\
What makes it viable
 & Program-level endorsement and visible employer demand
 & Advisor consent by default, and topics that feed the dissertation \\
\bottomrule
\end{tabularx}
\end{table}

We think the MS case is strong enough to act on now and the PhD case needs a different instrument.
For doctoral students, short projects are a poor substitute for the multi-month industrial placements
that funders already support, such as the NSF INTERN supplement. Their real value is earlier and
smaller: a Q1-sized project during the coursework year, before a student has committed to a narrow
research identity, is cheap enough that an advisor can say yes.

\section{Where: adoption in the Woodruff School}

The school's three program families do not face the same constraints, so a single rollout would stall
on the hardest one.

\begin{itemize}
  \item \textbf{ME} is the right starting point. Large MS enrollment, broad partner demand, and topics
        such as vibration analysis, thermal modeling, tolerance studies, or test-data reduction that are
        useful to a company and safe to hand to a student remotely.
  \item \textbf{NRE} needs a longer lead time. Export-control review under ITAR and EAR plus facility
        badging mean a brief cannot be approved in a week. Start with non-controlled work such as
        open-data analysis and code verification, and route controlled topics through existing
        sponsored-research agreements.
  \item \textbf{MP} must not collide with clinical training. CAMPEP accreditation and the residency
        pipeline already claim the practical hours, and patient data is bounded by HIPAA, so the opening
        is vendor-side and physics-support work rather than anything clinical.
\end{itemize}

We propose three phases. \textbf{Phase 0}, one term: an 8 to 10 week ME-only pilot with roughly 30 to
50 students, 5 to 8 companies, 8 to 12 mentors, and 12 to 20 briefs offering 40 to 60 actual places.
\textbf{Phase 1}, the next year: add NRE with export-control screening inside brief approval, and open
a doctoral track requiring advisor consent. \textbf{Phase 2}: add MP vendor projects and seek a
transcript notation or one-credit professional-practice course.

Phase 0 succeeds if at least 75\% of accepted engagements reach a reviewed completion, 90\% of
submissions get actionable feedback within five business days, and half the companies commit a second
brief. Those numbers test what actually breaks, which is student follow-through, mentor capacity, and
employer willingness to repeat.

\section{Case study: IndustryQuest}

To test feasibility rather than argue it, we built the operational layer as a working web application.
The system refuses a brief unless every target skill has a matching rubric criterion and at least one
criterion is marked as required to pass. Coordinators approve briefs, assign mentors, and own
escalation.

A representative engagement: a student completes a 2 to 6 hour vibration-analysis project for a pump
manufacturer. The mentor reviews the submission against the published rubric, requests one revision,
then scores the signal-analysis criterion 3 out of 4. The system records 120 points, a verified
completion, and \textit{Emerging} evidence in signal analysis. The account stays at level~1, because
levels track participation and are irrelevant to eligibility by design. That skill claim reaches
\textit{Bronze} only after a second project reviewed by a different assessor, so a student with more
points but one assessor is still not eligible for the harder project, and the interface says exactly
that instead of showing a match percentage.

The contribution is narrow and worth stating plainly. Brief approval, rubric-bound assessment, evidence
accumulation, and escalation can be run by one coordinator rather than one faculty member per team. It
does not show the learning outcomes are better. Only Phase 0 tests that.

\section{Pros, cons, and what would change our minds}

In favor: the model attacks coordination cost instead of adding a course, it produces assessment
evidence a program can report and an employer can check, and it degrades gracefully, since a student
who finishes one project has still done real reviewed work.

Against, and we take these seriously. Short projects cannot teach what a six-month immersion teaches,
and long ambiguity is missing by construction. Mentor capacity, not student demand, breaks first. A weak
rubric produces a confident and meaningless credential. Doctoral adoption needs advisor consent that
many advisors will reasonably withhold.

We would abandon the model if Phase 0 showed mentors unable to hold a five-day feedback turnaround at
even 40 to 60 places, or if two assessors scoring the same rubric criterion agreed no better than
chance. Both are measurable in one term, which is the point of starting small.

\begin{multicols}{2}\scriptsize
\begin{thebibliography}{99}
\setlength{\itemsep}{0pt}\setlength{\parskip}{0pt}
\bibitem{lave1991} J.~Lave and E.~Wenger, \textit{Situated Learning: Legitimate Peripheral Participation}. Cambridge University Press, 1991.
\bibitem{collins1989} A.~Collins, J.~S. Brown, and S.~E. Newman, ``Cognitive apprenticeship: Teaching the crafts of reading, writing, and mathematics,'' in \textit{Knowing, Learning, and Instruction}, L.~B. Resnick, Ed. Erlbaum, 1989, pp.~453--494.
\bibitem{nas2018} National Academies of Sciences, Engineering, and Medicine, \textit{Graduate STEM Education for the 21st Century}. National Academies Press, 2018.
\bibitem{crawley2014} E.~F. Crawley, J.~Malmqvist, S.~\"{O}stlund, D.~R. Brodeur, and K.~Edstr\"{o}m, \textit{Rethinking Engineering Education: The CDIO Approach}, 2nd ed. Springer, 2014.
\bibitem{wood1976} D.~Wood, J.~S. Bruner, and G.~Ross, ``The role of tutoring in problem solving,'' \textit{Journal of Child Psychology and Psychiatry}, vol.~17, no.~2, pp.~89--100, 1976.
\bibitem{biggs1996} J.~Biggs, ``Enhancing teaching through constructive alignment,'' \textit{Higher Education}, vol.~32, pp.~347--364, 1996.
\bibitem{gulikers2004} J.~T.~M. Gulikers, T.~J. Bastiaens, and P.~A. Kirschner, ``A five-dimensional framework for authentic assessment,'' \textit{Educational Technology Research and Development}, vol.~52, no.~3, pp.~67--86, 2004.
\bibitem{jonsson2007} A.~Jonsson and G.~Svingby, ``The use of scoring rubrics: Reliability, validity and educational consequences,'' \textit{Educational Research Review}, vol.~2, no.~2, pp.~130--144, 2007.
\bibitem{ericsson1993} K.~A. Ericsson, R.~T. Krampe, and C.~Tesch-R\"{o}mer, ``The role of deliberate practice in the acquisition of expert performance,'' \textit{Psychological Review}, vol.~100, no.~3, pp.~363--406, 1993.
\bibitem{kolb1984} D.~A. Kolb, \textit{Experiential Learning}. Prentice Hall, 1984.
\end{thebibliography}
\end{multicols}

\end{document}
